% Chapter 9 of Singular Learning Theory for the Rest of Us: the Bayes generalization error.
\sltchapter{Generalization}{The error $\lambda/n$ and the equations of state}{$y=ab^3x$}{ch:generalization}
\newtonfoot{2/6}{6}{\tfrac16}{a^2b^6}
\begin{refsection}[chapters/09-generalization.bib]

\begin{chaptersummary}
The regular theory says that the Bayes predictive distribution of a model with $d$ parameters has expected test loss $\frac d{2n}$ above the entropy of the truth. For the cubic cusp $y=ab^3x$ we compute by hand that the excess is half the squared posterior mean of the rescaled coefficient, a shrinkage estimator whose mean square is exactly $\frac13$: the excess is $\frac1{6n}$, not $\frac1n$. Two integrations by parts, a virial identity in the parameter and Stein's identity in the noise, prove $n(\E[G_n]-S)\to\lambda$ for every one-coefficient regression with Gaussian inputs and noise, and three of the four limits in Watanabe's equations of state for these models; his general theorem rests on the same two integrations by parts. Exact quadrature over $1000$ data sets per sample size confirms $\lambda$ for three models and, for the product model with $m=2$, shows how slowly the limit is reached. All of this concerns the exact posterior average at a fixed truth.
\end{chaptersummary}

\section{The question: what the posterior costs at test time}\label{sec:generalization:question}

A model trained on $n$ examples is judged on fresh ones, and when its output is a predictive distribution the natural score is the test cross-entropy, which no model can push below the entropy of the noise in the labels. How fast does the score approach that floor, and what sets the rate? For a Bayesian, who predicts with the posterior average of the models, the classical answer counts parameters. Whether that average is worth computing for neural networks is debated \parencite{WenzelEtAl2020,IzmailovEtAl2021}; this chapter asks what it buys in the simplest singular models, where it can be computed exactly.

Fix a model $p(x\mid w)$, $w\in W$, a prior $\varphi$, and data $X^n$ from the truth $q$. The \term{predictive distribution}\index{predictive distribution} is $p^*(x)=\Ew{p(x\mid w)}$, and its \term{Bayes generalization loss}\index{generalization loss} and Bayes training loss are
\[
G_n=-\E_X\bigl[\log p^*(X)\bigr],\qquad T_n=-\frac1n\sum_{i=1}^n\log p^*(X_i),
\]
with $X$ a fresh draw from $q$. If the model realizes the truth, $q=p(\cdot\mid w_0)$, the \term{generalization error}\index{generalization error} $G_n-S=\KL(q\,\|\,p^*)$ is the Kullback--Leibler divergence from the truth to the predictive; in a regression $p(y\mid x,w)$ the input density cancels, and $G_n-S$ averages over a test input the divergence between the conditional laws of $y$. The function $n\mapsto\E[G_n]-S$ is the \term{learning curve}\index{learning curve}.

\begin{regularanswer}[the generalization error]
If the truth is regular for a model with $d$ parameters and the model realizes it, then the Bayes predictive distribution has $\E[G_n]-S=\frac d{2n}+o(\frac1n)$ (\cref{thm:regular:aic}) and $\E[T_n]-S=-\frac d{2n}+o(\frac1n)$ \parencite{Watanabe2010a}; the maximum-likelihood plug-in has the same two values \parencite{AmariMurata1993}. Their difference $\frac dn$ is the penalty of Akaike's criterion on the per-sample scale \parencite{Akaike1974}.
\end{regularanswer}

The model of this chapter is the \term{cubic cusp}\index{cubic cusp} $y=ab^3x+\varepsilon$, with inputs $x\sim\Normal(0,1)$, unit noise, the uniform prior on $[-1,1]^2$, and the truth $(a,b)=(0,0)$, so that $y$ is pure noise. It has two parameters, and the box predicts an error of $\frac1n$.\sidenote{Akaike's criterion is usually written $-2\log L+2d$, with $L$ the maximized likelihood. That is $2n$ times $-\frac1n\log L+\frac dn$: the per-sample training loss plus the penalty $\frac dn$.} Its distributions $\Normal(cx,1)$, $c=ab^3$, are those of the regular line $y=cx$, $\abs c\le1$, which has one parameter and the error $\frac1{2n}$. Neither answer is right: the error is $\frac1{6n}$. Four models with the same distributions have three different limits (\cref{cor:generalization:running}): with the uniform prior on $[-1,1]^d$, the coefficient $c$ has the density $\frac12$ for the regular line $y=ax$, $\frac12\log\frac1{\abs c}$ for the product model $y=abx$, $\frac12(\abs c^{-1/2}-1)$ for the cusp $y=ab^2x$, and $\frac14(\abs c^{-2/3}-1)$ for the cubic cusp, and the errors are $\frac1{2n}$, $\frac1{2n}$, $\frac1{4n}$ and $\frac1{6n}$, against $\frac d{2n}=\frac1{2n},\frac1n,\frac1n,\frac1n$. What tells them apart is the prior that each induces on $c$ near the truth $c_0=0$, where the densities of the cusps blow up like a power of $\abs c$, and the product model's only like a logarithm: the parametrization bets on the truth, and the posterior average collects the bet. How does the generalization error of a singular model depend on $n$, and through which property of the model and the prior? The answer is $\E[G_n]-S=\lambda/n+o(1/n)$ at the temperature $\beta=1$ of Bayes' rule, with a companion number $\nu$ for the training error and for other temperatures.

\section{By hand: half the squared posterior mean}\label{sec:generalization:byhand}

A \term{one-coefficient model}\index{one-coefficient model} is a regression $y=g(w)x+\varepsilon$ whose mean depends on the parameters only through one number $c=g(w)$. Its posterior is the posterior of $c$ under the prior density $\rho$ that $\varphi$ induces on $c$ (\cref{ch:coefficient}), and its predictive at an input $x$ is a mixture of Gaussians, $p^*(y\mid x)=\Ew{\phi(y-cx)}$, with $\phi$ the standard normal density.

\begin{computation}{The divergence from the truth to the predictive}\label{comp:generalization:divergence}
\step Fix the data and a test input $x$; write $u=(c-c_0)x$ for the error of a coefficient drawn from the posterior and $\varepsilon=y-c_0x\sim\Normal(0,1)$ for the test noise. Since $\phi(y-cx)/\phi(y-c_0x)=e^{u\varepsilon-u^2/2}$,
\[
D_x=\KL\bigl(q(\cdot\mid x)\,\|\,p^*(\cdot\mid x)\bigr)=-\E_\varepsilon\log\Ew{e^{u\varepsilon-u^2/2}}.
\]
\step The probabilists' Hermite polynomials $\mathrm{He}_0=1$, $\mathrm{He}_1(\varepsilon)=\varepsilon$, $\mathrm{He}_2(\varepsilon)=\varepsilon^2-1$, \dots{} have the generating function $e^{u\varepsilon-u^2/2}=\sum_{k\ge0}\mathrm{He}_k(\varepsilon)u^k/k!$, so $\Ew{e^{u\varepsilon-u^2/2}}=1+A(\varepsilon)$ with
\[
A(\varepsilon)=\Ew{u}\varepsilon+\tfrac12\Ew{u^2}(\varepsilon^2-1)+\tfrac16\Ew{u^3}\mathrm{He}_3(\varepsilon)+\cdots.
\]
\step The polynomials are orthogonal under $\Normal(0,1)$, with $\E[\mathrm{He}_j\mathrm{He}_k]=k!$ if $j=k$ and $0$ otherwise. So $\E A=0$, and $\E A^2=\sum_{k\ge1}\Ew{u^k}^2/k!$.
\step Expanding $\log(1+A)=A-\frac12A^2+\frac13A^3-\cdots$, $D_x=\frac12\E A^2-\frac13\E A^3+\cdots=\frac12\Ew{u}^2+(\text{terms of degree }4\text{ in }u)$; terms of odd degree vanish because $\mathrm{He}_k(-\varepsilon)=(-1)^k\mathrm{He}_k(\varepsilon)$. The second moment $\Ew{u^2}$, which measures the spread of the posterior, multiplies $\mathrm{He}_2$, whose mean is $0$: the spread does not enter at second order.
\step The data locate $c$ to within about $n^{-1/2}$, so the terms of degree $4$ are of order $n^{-2}$. Averaging over the test input, with $\E x^2=1$,
\[
G_n-S=\E_xD_x=\tfrac12\bigl(\Ew{c}-c_0\bigr)^2+O(n^{-2}).
\]
\end{computation}

The generalization error is half the squared error of the posterior mean; a Gibbs predictor, which draws one coefficient from the posterior, pays $\frac12\Ew{(c-c_0)^2}$, the spread included. What made the computation possible was to read the likelihood ratio $e^{u\varepsilon-u^2/2}$ of a Gaussian shift as the generating function of an orthogonal family.

\begin{computation}{The regular line, exactly}\label{comp:generalization:line}
\step Take $y=ax$ with the flat prior on $\R$. With $S_{xx}=\sum_ix_i^2$ and $\hat c=\sum_ix_iy_i/S_{xx}$, the posterior of $c=a$ is $\Normal(\hat c,S_{xx}^{-1})$, and the predictive at $x$ is $\Normal(\hat cx,1+v)$ with $v=x^2/S_{xx}$, so
\[
D_x=\tfrac12\Bigl(\log(1+v)+\frac{1+(\hat c-c_0)^2x^2}{1+v}-1\Bigr).
\]
\step At the truth $Z_n=S_{xx}^{1/2}(\hat c-c_0)$ is standard normal and independent of $S_{xx}$ (\cref{lem:generalization:reduction}), and $(\hat c-c_0)^2x^2=Z_n^2v$. Averaging over $Z_n$ turns the fraction into exactly $1$, and $\E[D_x\mid S_{xx},x]=\frac12\log(1+v)$.
\step The squares $x^2\sim\chi^2_1$ and $S_{xx}\sim\chi^2_n$ are independent, so $B=x^2/(x^2+S_{xx})$ has the Beta law with parameters $\frac12$ and $\frac n2$, and $1+v=1/(1-B)$. With $\E\log(1-B)=\psi(\frac n2)-\psi(\frac{n+1}2)$, $\psi$ the digamma function,
\[
\E[G_n]-S=\tfrac12\Bigl(\psi\bigl(\tfrac{n+1}2\bigr)-\psi\bigl(\tfrac n2\bigr)\Bigr)=\frac1{2n}+\frac1{4n^2}+O(n^{-3}),
\]
the regular answer with $d=1$; at $n=100$, $n(\E[G_n]-S)=0.50250$. The rule of \cref{comp:generalization:divergence} gives $\frac12\E[(\hat c-c_0)^2]=\frac1{2(n-2)}$, which agrees to first order (\cref{ex:generalization:fourth}).
\end{computation}

\begin{computation}{Shrinkage in the cubic cusp}\label{comp:generalization:cusp}
\step For independent $a,b$ uniform on $[-1,1]$ and $0<s\le1$, $\Prob(\abs a\abs b^3\le s)=\int_0^1\min\bigl(1,(s/r)^{1/3}\bigr)\dd r=\frac32s^{1/3}-\frac12s$, so $c=ab^3$ has the density $\rho(c)=\frac14(\abs c^{-2/3}-1)$ on $[-1,1]$. Near $c_0=0$, $\rho(c)=\abs c^{-2/3}\ell(c)$ with $\ell(c)\to\frac14$, and the exponent $-\frac23=2\lambda-1$ gives the learning coefficient $\lambda=\frac16$ of $K=\frac12a^2b^6$.
\step Rescale the coefficient by the data, $t=S_{xx}^{1/2}c$. The energy is $nK_n=\frac12t^2-tZ_n$ (\cref{lem:generalization:reduction}), so the posterior of $t$ is proportional to $e^{-(t-Z_n)^2/2}(\abs t^{-2/3}-S_{xx}^{-1/3})$ on $\abs t\le S_{xx}^{1/2}$, and as $n\to\infty$ it tends to the density proportional to $e^{-(t-Z)^2/2}\abs t^{-2/3}$: a Gaussian likelihood centered at the least-squares coordinate $Z$, times a prior that is infinite at the truth. By \cref{comp:generalization:divergence}, $n(G_n-S)\approx\frac12\bar t^2$, with $\bar t$ the posterior mean of $t$.
\step For a Gaussian likelihood of unit variance and every prior weight $w$, \term{Tweedie's formula}\index{Tweedie's formula} \parencite{Robbins1956,Efron2011} gives the posterior mean $M(Z)=Z+h'(Z)/h(Z)$, with $h(Z)=\int e^{-(t-Z)^2/2}w(t)\dd t$ (\cref{thm:toolkit:stein}). For $w(t)=\abs t^{-2/3}$, expanding $e^{tZ}$ in $h(Z)=e^{-Z^2/2}\int\abs t^{-2/3}e^{-t^2/2+tZ}\dd t$, integrating term by term and applying Kummer's transformation give $h(Z)=2^{1/6}\Gamma(\frac16)\,{}_1F_1(\frac13;\frac12;-\frac{Z^2}2)$, with ${}_1F_1$ Kummer's confluent hypergeometric function, and therefore
\[
M(Z)=Z\Biggl(1-\frac23\,\frac{{}_1F_1\bigl(\frac43;\frac32;-\frac{Z^2}2\bigr)}{{}_1F_1\bigl(\frac13;\frac12;-\frac{Z^2}2\bigr)}\Biggr).
\]
\step The posterior pulls the least-squares coordinate toward the truth: $M(Z)\approx\frac13Z$ near $0$, where both functions ${}_1F_1$ equal $1$; $M(1)=0.412$, $M(2.5)=1.975$, $M(4)=3.804$; and $M(Z)=Z-\frac2{3Z}+O(Z^{-3})$ for large $\abs Z$. The posterior mean is a \term{shrinkage}\index{shrinkage} estimator of $Z$, strong near the truth and weak far from it.
\step Quadrature over $Z\sim\Normal(0,1)$ gives $\E[M(Z)^2]=0.333333\ldots$, which is $\frac13$ to ten digits. Hence $n(\E[G_n]-S)\to\frac16$: the cubic cusp pays $\frac1{6n}$, against $\frac1n$ for two regular parameters and $\frac1{2n}$ for the regular line.
\end{computation}

\begin{figure}[t]
\centering
\begin{tikzpicture}
\begin{axis}[splot,height=0.36\linewidth,xmin=0,xmax=5,ymin=0,ymax=2.75,ytick={0,0.5,1,1.5,2},
  xlabel={least-squares coordinate $Z$},ylabel={shrinkage factor $M(Z)/Z$},
  legend columns=2,legend style={font=\scriptsize,at={(0.98,0.97)},anchor=north east}]
\addplot[muted,line width=0.9pt,dashdotted] table[x=Z,y=S1] {figures/data/09-generalization-tweedie.dat};
\addlegendentry{$\lambda=1$}
\addplot[regularpred] table[x=Z,y=S12] {figures/data/09-generalization-tweedie.dat};
\addlegendentry{$\lambda=\frac12$ (regular)}
\addplot[theory2] table[x=Z,y=S14] {figures/data/09-generalization-tweedie.dat};
\addlegendentry{$\lambda=\frac14$ (cusp)}
\addplot[theory] table[x=Z,y=S16] {figures/data/09-generalization-tweedie.dat};
\addlegendentry{$\lambda=\frac16$ (cubic cusp)}
\end{axis}
\end{tikzpicture}
\caption{The shrinkage factor $M(Z)/Z$ of the mean of the limit posterior $\propto e^{-(t-Z)^2/2}\abs t^{2\lambda-1}$, by quadrature at $101$ values of $Z\ge0$ ($M$ is odd). It starts at $M'(0)=v(0)=2\lambda$ and tends to $1$: a prior weight that is infinite at the truth ($\lambda<\frac12$) shrinks the least-squares coordinate toward $0$, one that vanishes there ($\lambda=1$) pushes it away. In every case $\E[M(Z)^2]=2\lambda$ (\cref{lem:generalization:parts}). Script: \texttt{09-generalization-quartet.py}.}\label{fig:generalization:tweedie}
\end{figure}

Two things had to be noticed: at order $\frac1n$ only the posterior mean matters, and after the rescaling $t=S_{xx}^{1/2}(c-c_0)$ the problem forgets $n$, leaving a Gaussian location problem with one observation $Z$ and the prior weight $\abs t^{2\lambda-1}$ (\cref{fig:generalization:tweedie}). Why should the mean square of this shrinkage estimator be exactly $2\lambda$?\sidenote{Shrinkage estimators of a normal mean, and the identity used below to compute their risk, are the subject of \textcite{Stein1981}. Here the shrinkage is not designed: it is what the prior induced by the parametrization does, and it helps because the truth sits where that prior is singular.}

\section{The theorem: \texorpdfstring{$\lambda/n$}{lambda/n} for one-coefficient models}\label{sec:generalization:onecoef}

We prove the three facts the computation used, at every inverse temperature $\beta$, which \cref{sec:generalization:state} needs. At inverse temperature $\beta$ the posterior is proportional to $\varphi(w)e^{-n\beta L_n(w)}$, with expectations $\Ewb{\cdot}$, the predictive is $p^*_\beta(x)=\Ewb{p(x\mid w)}$, and $G_n(\beta)=-\E_X\log p^*_\beta(X)$, so that $G_n=G_n(1)$.

\begin{assumption}[One-coefficient regression]\label{ass:generalization:onecoef}
The data are independent pairs $(x_i,y_i)$, $i=1,\dots,n$, with $x_i\sim\Normal(0,1)$ and $y_i=c_0x_i+\varepsilon_i$, $\varepsilon_i\sim\Normal(0,1)$ independent of $x_i$. The model is $y=g(w)x+\varepsilon$, $\varepsilon\sim\Normal(0,1)$, for a measurable $g$ on $W\subseteq\R^d$ with $C_g=\sup_W\abs{g-c_0}<\infty$ and a prior probability density $\varphi$ on $W$; the truth is realizable, $c_0=g(w_0)$ for some $w_0\in W$. The law of $c=g(w)$ under $\varphi$ has a density $\rho$, and there are $\eta_0\in(0,1)$, $\lambda>0$, an integer $m\ge1$, and a function $\ell$ whose one-sided limits $\ell_\pm=\lim_{c\to c_0\pm}\ell(c)$ are both positive, such that
\[
\rho(c)=\abs{c-c_0}^{2\lambda-1}\Bigl(\log\frac1{\abs{c-c_0}}\Bigr)^{m-1}\ell(c)\qquad\text{for }0<\abs{c-c_0}<\eta_0.
\]
\end{assumption}

The exponent $\alpha=2\lambda-1>-1$ is fixed from now on, and $\rho$ vanishes outside $[c_0-C_g,c_0+C_g]$. The form of $\rho$ is that of \cref{thm:coefficient:free}, except that the density may approach different constants on the two sides of the truth. The numbers $\lambda$ and $m$ are the learning coefficient and multiplicity of $K(w)=\frac12(g(w)-c_0)^2$: the volume $\int_{\abs{c-c_0}<\sqrt{2\eps}}\rho(c)\dd c$ is a constant times $\eps^\lambda(\log\frac1\eps)^{m-1}(1+o(1))$ (\cref{ch:volume}).\sidenote{The learning coefficient is the real log canonical threshold of $K$ with respect to $\varphi$: two names for one number.}

\begin{lemma}[The data enter through two numbers]\label{lem:generalization:reduction}
Under \cref{ass:generalization:onecoef} let $S_{xx}=\sum_ix_i^2$, $\kappa_n=S_{xx}^{1/2}$ and $Z_n=\kappa_n^{-1}\sum_ix_i\varepsilon_i$, and for $w\in W$ let $t=\kappa_n(g(w)-c_0)$.
\begin{enumerate}[label=(\roman*),nosep]
\item The variables $\kappa_n^2\sim\chi^2_n$ and $Z_n\sim\Normal(0,1)$ are independent.
\item $nK_n(w)=\frac12t^2-tZ_n$ for every $w$.
\item At inverse temperature $\beta$ the posterior law of $t$ has the density proportional to $e^{-\beta(t-Z_n)^2/2}\rho(c_0+t/\kappa_n)$.
\item $G_n(\beta)-S=\E_x\bigl[D(\mu_x)\bigr]$, where $x\sim\Normal(0,1)$, $\mu_x$ is the posterior law of $u=(g(w)-c_0)x$, and for a probability distribution $\mu$ on $\R$ and $\varepsilon\sim\Normal(0,1)$
\[
D(\mu)=-\E_\varepsilon\log\int e^{u\varepsilon-u^2/2}\mu(\dd u).
\]
\end{enumerate}
\end{lemma}

\begin{proof}
The plan is to condition on the inputs and expand the squares. (i) Given the inputs, $\sum_ix_i\varepsilon_i$ is Gaussian with mean $0$ and variance $S_{xx}$, so $Z_n$ is standard normal given the inputs; its conditional law does not depend on them, which makes $Z_n$ independent of the inputs, and of $\kappa_n$. (ii) With $c=g(w)$, $y_i-cx_i=\varepsilon_i-(c-c_0)x_i$; the normalizing constants cancel in $nK_n=nL_n(w)-nL_n(w_0)$, and
\[
nK_n(w)=\tfrac12\sum_i\bigl((\varepsilon_i-(c-c_0)x_i)^2-\varepsilon_i^2\bigr)=\tfrac12t^2-tZ_n .
\]
(iii) The posterior is proportional to $\varphi(w)e^{-n\beta K_n(w)}$, which depends on $w$ only through $c$. For bounded $F$ the definition of $\rho$ gives, with $t=\kappa_n(c-c_0)$,
\[
\int F(g(w))\,e^{-n\beta K_n(w)}\varphi(w)\dd w=\int F(c)\,e^{-\beta(t^2/2-tZ_n)}\rho(c)\dd c,
\]
and completing the square gives (iii). (iv) The input density cancels, and for a test pair with $\varepsilon=y-c_0x$, $\phi(y-cx)/\phi(y-c_0x)=e^{u\varepsilon-u^2/2}$, so $p^*_\beta(y\mid x)/q(y\mid x)=\int e^{u\varepsilon-u^2/2}\mu_x(\dd u)$, with $\varepsilon\sim\Normal(0,1)$ under the truth.
\end{proof}

The function $D(\mu)$ is the Kullback--Leibler divergence from $\Normal(0,1)$ to the mixture $\int\Normal(u,1)\,\mu(\dd u)$, and \cref{comp:generalization:divergence} becomes the following inequality.

\begin{lemma}[The divergence to a Gaussian mixture]\label{lem:generalization:divergence}
Let $\mu$ be a probability distribution on $[-a,a]$, $a>0$, with mean $\bar u$ and second moment $\gamma=\int u^2\mu(\dd u)$. Then
\begin{enumerate}[label=(\roman*),nosep]
\item $0\le D(\mu)\le\frac12\gamma$, and hence $\bigl\lvert D(\mu)-\frac12\bar u^2\bigr\rvert\le\frac12\gamma$;
\item $\bigl\lvert D(\mu)-\frac12\bar u^2\bigr\rvert\le\frac14(1+a^2)^2e^{9a^2}\gamma^2$.
\end{enumerate}
\end{lemma}

\begin{proof}
The plan is to split off the Gaussian at the mean, $\Normal(\bar u,1)$, whose divergence from $\Normal(0,1)$ is exactly $\frac12\bar u^2$, and to bound the rest with a quadratic bound for the logarithm, after Jensen's inequality has kept the mixture away from $0$.

(i) The divergence is nonnegative. By Jensen's inequality for the logarithm, $-\log\int e^{u\varepsilon-u^2/2}\dd\mu\le\int(\frac12u^2-u\varepsilon)\dd\mu$, whose mean over $\varepsilon$ is $\frac12\gamma$; and $0\le\frac12\bar u^2\le\frac12\gamma$.

(ii) Write $r=u-\bar u$, so that $\int r\dd\mu=0$, $\gamma_0=\int r^2\dd\mu=\gamma-\bar u^2$, $\abs r\le2a$ and $\abs{\bar u}\le a$. Since $u\varepsilon-\frac12u^2=(\bar u\varepsilon-\frac12\bar u^2)+(r(\varepsilon-\bar u)-\frac12r^2)$,
\[
D(\mu)=\tfrac12\bar u^2+T,\qquad T=-\E_\varepsilon\log N(\varepsilon-\bar u),\qquad N(z)=\int e^{rz-r^2/2}\mu(\dd u).
\]
Jensen's inequality and $\int r\dd\mu=0$ give $N(z)\ge e^{-\gamma_0/2}$, and $N(z)\le e^{2a\abs z}$, so $T$ is finite. We use $\abs{e^y-1}\le\abs ye^{\abs y}$, $\abs{e^y-1-y}\le\frac12y^2e^{\abs y}$, $\abs{r\bar u}\le2a^2$, and $\abs{rr'}\le4a^2$ for an independent copy $r'$ of $r$.

\emph{Lower bound.} Since $-\log N\ge1-N$ and $\E_\varepsilon e^{r(\varepsilon-\bar u)-r^2/2}=e^{-r\bar u}$ (Tonelli's theorem exchanges the integrals),
\[
T\ge-\int\bigl(e^{-r\bar u}-1+r\bar u\bigr)\dd\mu\ge-\tfrac12e^{2a^2}\bar u^2\gamma_0\ge-\tfrac18e^{2a^2}\gamma^2,
\]
because $\bar u^2\gamma_0\le\frac14(\bar u^2+\gamma_0)^2=\frac14\gamma^2$.

\emph{Upper bound.} Taylor's theorem gives $-\log y=(1-y)+\frac{(1-y)^2}{2\xi^2}$ with $\xi$ between $y>0$ and $1$; for $y=N(z)$, $\xi\ge e^{-\gamma_0/2}$, so $-\log N\le(1-N)+\frac12e^{\gamma_0}(1-N)^2$. The mean of $1-N(\varepsilon-\bar u)$ is $1-\int e^{-r\bar u}\dd\mu\le0$ by Jensen's inequality, so $T\le\frac12e^{\gamma_0}\E_\varepsilon(N(\varepsilon-\bar u)-1)^2$, and expanding the square with $\E_\varepsilon e^{(r+r')(\varepsilon-\bar u)-(r^2+r'^2)/2}=e^{rr'-(r+r')\bar u}$,
\[
\E_\varepsilon(N-1)^2=\iint\bigl(e^{rr'}-1\bigr)e^{-(r+r')\bar u}\dd\mu\dd\mu+\Bigl(\int e^{-r\bar u}\dd\mu-1\Bigr)^2 .
\]
The part of the double integral with $rr'$ in place of $e^{rr'}-1$ is $\bigl(\int r(e^{-r\bar u}-1)\dd\mu\bigr)^2\le e^{4a^2}\bar u^2\gamma_0^2$; the rest of it is at most $\frac12\iint(rr')^2e^{\abs{rr'}+\abs{r\bar u}+\abs{r'\bar u}}\le\frac12e^{8a^2}\gamma_0^2$; the last square is at most $\frac14e^{4a^2}\bar u^4\gamma_0^2$. With $\bar u^2\le a^2$ the sum is at most $\gamma_0^2e^{8a^2}(\frac12+a^2+\frac14a^4)\le\frac12(1+a^2)^2e^{8a^2}\gamma_0^2$, and since $\gamma_0\le\gamma\le a^2$, $T\le\frac14(1+a^2)^2e^{9a^2}\gamma^2$.
\end{proof}

The next lemma controls the rescaled posterior of \cref{lem:generalization:reduction}(iii) uniformly in the data and identifies its limit.

\begin{lemma}[The rescaled posterior]\label{lem:generalization:posterior}
Under \cref{ass:generalization:onecoef} let $\beta>0$. For $\kappa>1$ and $z\in\R$ let $\pi_{\kappa,z}$ be the probability density proportional to $e^{-\beta(t-z)^2/2}\rho(c_0+t/\kappa)$, and $\pi_z$ the one proportional to $e^{-\beta(t-z)^2/2}w(t)$, with $w(t)=\ell_+t^\alpha$ for $t>0$ and $w(t)=\ell_-\abs t^\alpha$ for $t<0$. There are constants $\kappa_1$ and $C$ such that:
\begin{enumerate}[label=(\roman*),nosep]
\item $\int t^2\pi_z(t)\dd t\le C(1+z^2)$ for every $z$, a bound that uses only $\alpha>-1$ and $\ell_\pm>0$; and $\int t^2\pi_{\kappa,z}(t)\dd t\le C(1+z^2)$ whenever $\kappa\ge\kappa_1$ and $\abs z\le\sqrt\kappa-1$;
\item for every $z$ and $k\in\{1,2\}$, $\int t^k\pi_{\kappa,z}(t)\dd t\to\int t^k\pi_z(t)\dd t$ as $\kappa\to\infty$.
\end{enumerate}
\end{lemma}

\begin{proof}
The plan is to rescale the prior density so that it converges to $w$, to bound it above and below on the scales that matter, and to compare the numerators and denominators of the posterior moments.

\emph{The rescaled density.} Shrinking $\eta_0$, we may assume $\frac12\ell_\pm\le\ell(c)\le2\ell_\pm$ for $0<\pm(c-c_0)<\eta_0$; put $\bar\ell=2\max\ell_\pm$, $\underline\ell=\frac12\min\ell_\pm$, and $\rho_\kappa(t)=\kappa^\alpha(\log\kappa)^{1-m}\rho(c_0+t/\kappa)$, so that $\pi_{\kappa,z}$ is proportional to $e^{-\beta(t-z)^2/2}\rho_\kappa(t)$. For $0<\abs t<\eta_0\kappa$,
\[
\rho_\kappa(t)=\abs t^\alpha\Bigl(1-\frac{\log\abs t}{\log\kappa}\Bigr)^{m-1}\ell(c_0+t/\kappa).
\]
Let $\kappa\ge e$. The bracket lies between $1$ and $1+\log\frac1{\abs t}$ if $\abs t\le1$, in $(0,1]$ if $1\le\abs t<\eta_0\kappa$, and is at least $\frac12$ if $\abs t\le\sqrt\kappa$. Hence, with $\log^+y=\max(\log y,0)$,
\begin{gather*}
\rho_\kappa(t)\le\bar\ell\,\abs t^\alpha\bigl(1+\log^+\tfrac1{\abs t}\bigr)^{m-1}\qquad(0<\abs t<\eta_0\kappa),\\
\rho_\kappa(t)\ge2^{1-m}\underline\ell\,\abs t^\alpha\qquad(0<\abs t\le\sqrt\kappa<\eta_0\kappa),\qquad
\int_{\abs t\ge\eta_0\kappa}\rho_\kappa(t)\dd t\le\kappa^{1+\alpha}.
\end{gather*}
For each $t\ne0$, $\rho_\kappa(t)\to w(t)$ as $\kappa\to\infty$; and $w$ satisfies the two pointwise bounds for every $t\ne0$.

\emph{Denominators.} Let $\abs z\le\sqrt\kappa-1$, let $\varsigma$ be $1$ if $z\ge0$ and $-1$ otherwise, and let $J$ be the interval between $z+\frac12\varsigma$ and $z+\varsigma$. On $J$, $\frac12(1+\abs z)\le\abs t\le1+\abs z\le\sqrt\kappa$ and $e^{-\beta(t-z)^2/2}\ge e^{-\beta/2}$, so
\[
\int e^{-\beta(t-z)^2/2}\rho_\kappa(t)\dd t\ge c_1(1+\abs z)^\alpha,\qquad c_1=\tfrac12e^{-\beta/2}2^{1-m-\abs\alpha}\underline\ell,
\]
and the same holds with $w$ in place of $\rho_\kappa$ for every $z$.

\emph{Numerators.} Let $\kappa_1\ge e$ be so large that, for $\kappa\ge\kappa_1$, $\sqrt\kappa\le\frac12\eta_0\kappa$ and $\kappa^{1+\alpha}s^2e^{-\beta s^2/8}\le1$ for every $s\ge\eta_0\kappa$. On $\abs t<\eta_0\kappa$ the upper bound reduces the integral of $t^2e^{-\beta(t-z)^2/2}\rho_\kappa(t)$ to $\bar\ell I(z)$, where, splitting at $\abs t=1$,
\begin{align*}
I(z)&=\int\abs t^{2+\alpha}\bigl(1+\log^+\tfrac1{\abs t}\bigr)^{m-1}e^{-\beta(t-z)^2/2}\dd t\\
&\le\int_{-1}^1\abs t^\alpha\bigl(1+\log\tfrac1{\abs t}\bigr)^{m-1}\dd t+\int\abs t^{2+\alpha}e^{-\beta(t-z)^2/2}\dd t\le C_2(1+\abs z)^{2+\alpha}
\end{align*}
for a constant $C_2$: the first integral is finite because $\alpha>-1$, and in the second $\abs t^{2+\alpha}\le2^{1+\alpha}(\abs z^{2+\alpha}+\abs{t-z}^{2+\alpha})$ by convexity, as $2+\alpha\ge1$. On $\abs t\ge\eta_0\kappa$, $\abs{t-z}\ge\frac12\abs t$ because $\abs z\le\frac12\eta_0\kappa$, and the integral is at most $1$.

\emph{Conclusion of (i).} Dividing, $\int t^2\pi_{\kappa,z}\le\bigl(\bar\ell C_2(1+\abs z)^{2+\alpha}+1\bigr)/\bigl(c_1(1+\abs z)^\alpha\bigr)\le C(1+z^2)$, since $(1+\abs z)^{-\alpha}\le(1+\abs z)^2\le2(1+z^2)$. The same computation without the far part bounds the moments of $\pi_z$.

\emph{Convergence.} Fix $z$; for large $\kappa$, $\abs z\le\frac12\eta_0\kappa$. For $k=0,1,2$ the integrand $t^ke^{-\beta(t-z)^2/2}\rho_\kappa(t)\ind\{\abs t<\eta_0\kappa\}$ converges to $t^ke^{-\beta(t-z)^2/2}w(t)$ for $t\ne0$ and is dominated by the integrable $\bar\ell\abs t^{k+\alpha}(1+\log^+\frac1{\abs t})^{m-1}e^{-\beta(t-z)^2/2}$, while the far part tends to $0$. By dominated convergence the three moment integrals converge, the zeroth to a positive limit, and so do the ratios.
\end{proof}

The weight of the limit posterior is homogeneous of degree $\alpha$, $w(st)=s^\alpha w(t)$ for $s>0$. Two integrations by parts compute the mean square of the posterior mean: one in the parameter, which uses the homogeneity and which we call the \term{virial identity}\index{virial identity}, after its counterpart in mechanics;\sidenote{The virial theorem averages $x\cdot\nabla H$ against a Gibbs measure. Here $H=-\log\pi_z$, and the average of $t\,\partial_tH=\beta t(t-z)-\alpha$ is $1$, by one integration by parts: that is (i) of \cref{lem:generalization:parts}.} and one in the noise, \term{Stein's identity}\index{Stein's identity} $\E[Zf(Z)]=\E[f'(Z)]$ for $Z\sim\Normal(0,1)$ and absolutely continuous $f$ with $\E\abs{f'(Z)}<\infty$ (\cref{thm:toolkit:stein}; \textcite{Stein1981}).

\begin{lemma}[Two integrations by parts]\label{lem:generalization:parts}
Let $\alpha>-1$, $\ell_\pm>0$, $\beta>0$, and let $\pi_z$ and $w$ be as in \cref{lem:generalization:posterior}, with mean $M(z)$ and variance $v(z)$. Then:
\begin{enumerate}[label=(\roman*),nosep]
\item (virial identity) $\beta\int t(t-z)\pi_z(t)\dd t=1+\alpha$ for every $z\in\R$;
\item (Tweedie) $M$ is continuously differentiable and $M'(z)=\beta v(z)$;
\item (Stein) for $Z\sim\Normal(0,1)$, $\E[ZM(Z)]=\beta\E[v(Z)]$, all terms being finite;
\item $\E[M(Z)^2]=\frac{1+\alpha}\beta-(1-\beta)\E[v(Z)]$; in particular $\E[M(Z)^2]=1+\alpha=2\lambda$ at $\beta=1$.
\end{enumerate}
\end{lemma}

\begin{proof}
The plan is to integrate by parts in $t$ for (i), to differentiate under the integral sign for (ii), and to combine the two with Stein's identity. Write $\psi_z(t)=e^{-\beta(t-z)^2/2}$ and $h_k(z)=\int t^k\psi_z(t)w(t)\dd t$, finite for $k\ge0$ because $\alpha>-1$; then $M=h_1/h_0$ and $v=h_2/h_0-M^2$.

(i) On $(0,\infty)$, $\beta(t-z)\psi_z=-\partial_t\psi_z$, and $tw(t)=\ell_+t^{1+\alpha}$ tends to $0$ as $t\downarrow0$, because $1+\alpha>0$, and has the derivative $(1+\alpha)w(t)$. Integrating by parts, the boundary terms vanish and
\[
\beta\int_0^\infty t(t-z)\psi_z(t)w(t)\dd t=(1+\alpha)\int_0^\infty\psi_z(t)w(t)\dd t .
\]
On $(-\infty,0)$, $tw(t)=-\ell_-\abs t^{1+\alpha}$ also has the derivative $(1+\alpha)w(t)$, and the same computation applies. Add the two halves and divide by $h_0(z)$.

(ii) For $\abs z\le A$, $\abs{t-z}\psi_z(t)\le C_Ae^{-\beta t^2/8}$, because $(t-z)^2\ge\frac12t^2-A^2$; so the $h_k$ may be differentiated under the integral sign, with $h_k'=\beta(h_{k+1}-zh_k)$, and $M'=(h_1'h_0-h_1h_0')/h_0^2=\beta(h_2/h_0-h_1^2/h_0^2)=\beta v$, which is continuous. At $\beta=1$, $h_0'/h_0=M-z$ is Tweedie's formula.

(iii) By \cref{lem:generalization:posterior}(i), $M(z)^2\le\int t^2\pi_z\le C(1+z^2)$ and $0\le v(z)\le C(1+z^2)$. So $\E\abs{M'(Z)}=\beta\E v(Z)<\infty$, and Stein's identity gives $\E[ZM(Z)]=\E[M'(Z)]=\beta\E[v(Z)]$.

(iv) Since $\int t^2\pi_z=M^2+v$, (i) says $M(z)^2+v(z)-zM(z)=(1+\alpha)/\beta$ for every $z$. Put $z=Z$ and take expectations, which are finite by (iii); then (iii) gives
\[
\E[M(Z)^2]=\frac{1+\alpha}\beta-\E[v(Z)]+\E[ZM(Z)]=\frac{1+\alpha}\beta-(1-\beta)\E[v(Z)] .\qedhere
\]
\end{proof}

For the symmetric weight $\abs t^\alpha$ at $\beta=1$, the virial identity is the heat-equation identity $h''=\alpha h-zh'$ of \cref{thm:toolkit:stein}(iv) in other words: with Tweedie's $M=z+r$, $r=h'/h$, both say $r'=\alpha-zr-r^2$ (\cref{ex:generalization:riccati}). The virial form needs neither the symmetry nor $\beta=1$.

\begin{theorem}[$\lambda/n$ for one-coefficient models]\label{thm:generalization:onecoef}
Under \cref{ass:generalization:onecoef}, at $\beta=1$:
\begin{enumerate}[label=(\roman*),nosep]
\item $n\bigl(\E[G_n]-S\bigr)\to\lambda$ as $n\to\infty$, the expectation being over the data;
\item $n(G_n-S)$ converges in law to $\frac12M(Z)^2$, with $Z\sim\Normal(0,1)$ and $M$ the mean of the limit posterior $\pi_z$ of \cref{lem:generalization:posterior} at $\beta=1$.
\end{enumerate}
\end{theorem}

\begin{proof}
The plan is to replace the divergence by half the squared posterior mean (\cref{lem:generalization:divergence}), to control the posterior on an event of overwhelming probability (\cref{lem:generalization:posterior}), and to pass to the limit, where \cref{lem:generalization:parts} evaluates the mean square. The first three steps hold at every $\beta>0$, and we prove them so, for \cref{prop:generalization:quartet}. Write $\pi_n=\pi_{\kappa_n,Z_n}$ for the posterior of $t=\kappa_n(c-c_0)$ (\cref{lem:generalization:reduction}(iii)), $\bar t_n$ and $\Theta_n$ for its mean and second moment, and $M=M_\beta$.

\emph{From the divergence to the posterior mean.} For a test input $x$, the posterior law $\mu_x$ of $u=tx/\kappa_n$ has mean $x\bar t_n/\kappa_n$ and second moment $x^2\Theta_n/\kappa_n^2$, and lives on $\abs u\le C_g\abs x$. Put $\Delta_x=\bigl\lvert D(\mu_x)-x^2\bar t_n^2/(2\kappa_n^2)\bigr\rvert$. \Cref{lem:generalization:divergence}(i) gives $\Delta_x\le x^2\Theta_n/(2\kappa_n^2)$ for every $x$, and part (ii) with $a=C_gK$ gives, for $\abs x\le K$,
\[
\Delta_x\le C_K\,\frac{x^4\Theta_n^2}{\kappa_n^4},\qquad C_K=\tfrac14(1+C_g^2K^2)^2e^{9C_g^2K^2}.
\]
Average over $x\sim\Normal(0,1)$, with $\E x^2=1$, $\E x^4=3$ and $e_K=\E[x^2\ind\{\abs x>K\}]$, and use \cref{lem:generalization:reduction}(iv):
\begin{equation}\label{eq:generalization:step1}
\Bigl\lvert n\bigl(G_n(\beta)-S\bigr)-\frac{n\bar t_n^2}{2\kappa_n^2}\Bigr\rvert\le3C_K\frac{n\Theta_n^2}{\kappa_n^4}+\frac{e_K}2\,\frac{n\Theta_n}{\kappa_n^2}.
\end{equation}

\emph{A good event.} Let $A_n$ be the event that $\kappa_n^2\ge n/2$ and $\abs{Z_n}\le(n/2)^{1/4}-1$. Chernoff's bound gives $\Prob(\chi^2_n<\frac n2)\le e^{n/4}\E e^{-\chi^2_n/2}=e^{n/4}2^{-n/2}\le e^{-n/11}$, and $\Prob(\abs{Z_n}>s)\le2e^{-s^2/2}$; so $n^2\Prob(A_n^c)\to0$. Let $n$ be so large that $\sqrt{n/2}\ge\kappa_1$. On $A_n$, $\kappa_n\ge\kappa_1$ and $\abs{Z_n}\le\sqrt{\kappa_n}-1$, so $\Theta_n\le C(1+Z_n^2)$ by \cref{lem:generalization:posterior}(i), while $n/\kappa_n^2\le2$; on every data set $\abs t\le C_g\kappa_n$ under $\pi_n$, so $\Theta_n\le C_g^2\kappa_n^2$. Therefore
\begin{gather*}
\E\Bigl[\frac{n\Theta_n}{\kappa_n^2}\Bigr]\le2C\,\E[1+Z^2]+nC_g^2\Prob(A_n^c)\to4C,\\
\E\Bigl[\frac{n\Theta_n^2}{\kappa_n^4}\Bigr]\le\frac{4C^2}n\E\bigl[(1+Z^2)^2\bigr]+nC_g^4\Prob(A_n^c)\to0 .
\end{gather*}
Taking expectations in \eqref{eq:generalization:step1}, $\limsup_n\bigl\lvert n\E[G_n(\beta)-S]-\frac12\E[n\bar t_n^2/\kappa_n^2]\bigr\rvert\le2Ce_K$ for every $K$, and $e_K\to0$ as $K\to\infty$.

\emph{The posterior mean converges.} Let $\bar t(\kappa,z)$ be the mean of $\pi_{\kappa,z}$, a deterministic function, and $\Psi(\kappa)=\int_{\abs z\le\sqrt\kappa-1}\bar t(\kappa,z)^2\phi(z)\dd z$. For $\kappa\ge\kappa_1$ the integrand is at most $C(1+z^2)\phi(z)$ by \cref{lem:generalization:posterior}(i), and it converges to $M(z)^2\phi(z)$ by part (ii); so $\Psi\le2C$ and $\Psi(\kappa)\to\E[M(Z)^2]$ by dominated convergence. Let $B_n\supseteq A_n$ be the event that $\kappa_n^2\ge n/2$ and $\abs{Z_n}\le\sqrt{\kappa_n}-1$. Since $Z_n\sim\Normal(0,1)$ independently of $\kappa_n$ (\cref{lem:generalization:reduction}(i)),
\[
\E\Bigl[\frac{n\bar t_n^2}{\kappa_n^2}\ind_{B_n}\Bigr]=\E\Bigl[\frac n{\kappa_n^2}\ind\{\kappa_n^2\ge n/2\}\,\Psi(\kappa_n)\Bigr].
\]
The variable on the right is at most $4C$ and converges in probability to $\E[M(Z)^2]$, because $\kappa_n^2/n\to1$ in probability by the law of large numbers; by bounded convergence the right side tends to $\E[M(Z)^2]$. Off $B_n$, $\bar t_n^2\le\Theta_n\le C_g^2\kappa_n^2$, and $nC_g^2\Prob(B_n^c)\to0$. With the previous step,
\begin{equation}\label{eq:generalization:anybeta}
n\bigl(\E[G_n(\beta)]-S\bigr)\to\tfrac12\E\bigl[M_\beta(Z)^2\bigr]\qquad\text{for every }\beta>0.
\end{equation}

\emph{Conclusion.} At $\beta=1$, \cref{lem:generalization:parts}(iv) gives $\frac12\E[M(Z)^2]=\frac12(1+\alpha)=\lambda$, which is (i). For (ii), the expectation of the left side of \eqref{eq:generalization:step1} has limit superior at most $2Ce_K$ for every $K$, so it tends to $0$; with $n/\kappa_n^2\to1$ in probability and $\bar t_n^2\le\Theta_n=O_p(1)$, $n(G_n-S)-\frac12\bar t_n^2\to0$ in probability. The function $\Xi(\kappa)=\int\min(1,\abs{\bar t(\kappa,z)-M(z)})\phi(z)\dd z$ tends to $0$ by \cref{lem:generalization:posterior}(ii) and bounded convergence, and by independence $\E\min(1,\abs{\bar t_n-M(Z_n)})=\E\,\Xi(\kappa_n)\to0$. So $n(G_n-S)-\frac12M(Z_n)^2\to0$ in probability, and $\frac12M(Z_n)^2$ has the law of $\frac12M(Z)^2$ for every $n$; Slutsky's lemma (\cref{def:toolkit:convergence}) gives (ii).
\end{proof}

\begin{corollary}[The running models]\label{cor:generalization:running}
With inputs $\Normal(0,1)$, unit noise and the uniform prior on $[-1,1]^d$, at the truth $c_0=0$, $n(\E[G_n]-S)$ tends to $\frac12$ for the regular line $y=ax$, the product model $y=abx$ and the chain $y=a_1\cdots a_Lx$, to $\frac14$ for the cusp $y=ab^2x$, and to $\frac16$ for the cubic cusp $y=ab^3x$. At a truth with $0<\abs{c_0}<1$ it tends to $\frac12$ for all five.
\end{corollary}

\begin{proof}
The densities of $c$, given in \cref{sec:generalization:question} and \cref{comp:generalization:cusp} and equal to $(\log\frac1{\abs c})^{L-1}/(2(L-1)!)$ for the chain (\cref{ch:coefficient}), have the form of \cref{ass:generalization:onecoef} at $c_0=0$ with $(\lambda,m)=(\frac12,1),(\frac12,2),(\frac12,L),(\frac14,1),(\frac16,1)$ and $\ell(0\pm)>0$; they are continuous and positive at every $c_0$ with $0<\abs{c_0}<1$, the form with $\lambda=\frac12$, $m=1$; and $\abs c\le1$.
\end{proof}

The multiplicity entered only through the factor $(1-\log\abs t/\log\kappa)^{m-1}$, which tends to $1$: the limit is $\lambda$ whatever $m$ is, but for $m\ge2$ the approach is slow (\cref{sec:generalization:measuring}). Gaussianity entered through the exact law of $\kappa_n$ and the independence of $Z_n$ and $\kappa_n$. The proof is the book's own, the one-coefficient case of the two integrations by parts in Watanabe's general proof (\cref{sec:generalization:state}).\sidenote{We have not found the argument in the literature. The chapter's script confirms $\E[M(Z)^2]=1+\alpha$ to ten digits for $\lambda=\frac16,\frac14,\frac12,\frac34,1$ and for asymmetric weights.}

\section{The theorem: the equations of state}\label{sec:generalization:state}

A Bayesian can score the posterior average (Bayes) or a single posterior draw (\term{Gibbs}\index{Gibbs error}), on test or training data: at inverse temperature $\beta$, besides $G_n(\beta)$, there are the Bayes training loss $T_n(\beta)=-\frac1n\sum_i\log p^*_\beta(X_i)$ and the Gibbs losses
\[
G_{n,\mathrm g}(\beta)=\Ewb{L(w)},\qquad T_{n,\mathrm g}(\beta)=\Ewb{L_n(w)} .
\]
In the realizable case the four errors are $G_n(\beta)-S$, $T_n(\beta)-S_n$, $G_{n,\mathrm g}(\beta)-S=\Ewb{K(w)}$ and $T_{n,\mathrm g}(\beta)-S_n=\Ewb{K_n(w)}$; only the training losses can be computed from the data alone. The \term{functional variance}\index{functional variance} is $V_n=\sum_i\Var^\beta_w[\log p(X_i\mid w)]$.

\begin{proposition}[The four errors of a one-coefficient model]\label{prop:generalization:quartet}
Under \cref{ass:generalization:onecoef}, let $\beta>0$, let $M$ and $v$ be the mean and variance of the limit posterior $\pi_z$ of \cref{lem:generalization:posterior}, and let $\nu(\beta)=\frac\beta2\E[v(Z)]$ with $Z\sim\Normal(0,1)$. As $n\to\infty$,
\begin{gather*}
n\E\bigl[G_n(\beta)-S\bigr]\to\frac{\lambda-\nu}\beta+\nu,\\
n\E\bigl[G_{n,\mathrm g}(\beta)-S\bigr]\to\frac\lambda\beta+\nu,\qquad n\E\bigl[T_{n,\mathrm g}(\beta)-S_n\bigr]\to\frac\lambda\beta-\nu .
\end{gather*}
\end{proposition}

\begin{proof}
The plan is to read the three limits off the proof of \cref{thm:generalization:onecoef} and to evaluate them with \cref{lem:generalization:parts}. By \eqref{eq:generalization:anybeta} the first limit is $\frac12\E[M(Z)^2]$, and \cref{lem:generalization:parts}(iv) with $\E v=2\nu/\beta$ turns it into $\frac{1+\alpha}{2\beta}-\frac{(1-\beta)\nu}\beta=\frac{\lambda-\nu}\beta+\nu$. For the Gibbs errors, $K(w)=\frac12(c-c_0)^2\E x^2=t^2/(2\kappa_n^2)$ and $nK_n=\frac12t^2-tZ_n$ (\cref{lem:generalization:reduction}), so
\[
n\bigl(G_{n,\mathrm g}(\beta)-S\bigr)=\frac{n\Theta_n}{2\kappa_n^2},\qquad n\bigl(T_{n,\mathrm g}(\beta)-S_n\bigr)=\tfrac12\Theta_n-Z_n\bar t_n .
\]
The step ``the posterior mean converges'' of that proof, with $\int t^2\pi_{\kappa,z}$ and $z\,\bar t(\kappa,z)$ in place of $\bar t(\kappa,z)^2$, both at most a constant times $1+z^2$ for $\kappa\ge\kappa_1$ and $\abs z\le\sqrt\kappa-1$, shows that the expectations on the good event converge to $\frac12\E[M^2+v]$ and $\frac12\E[M^2+v]-\E[ZM]$. Off it, $n\Theta_n/\kappa_n^2\le nC_g^2$, $\Theta_n\le C_g^2\kappa_n^2$ and $\abs{Z_n\bar t_n}\le C_g\abs{Z_n}\kappa_n$, whose expectations on $A_n^c$ are at most $nC_g^2\Prob(A_n^c)$, $C_g^2(\E\kappa_n^4)^{1/2}\Prob(A_n^c)^{1/2}$ and $C_g\sqrt n\,\Prob(A_n^c)^{1/2}$ by the Cauchy--Schwarz inequality; they tend to $0$ because $\E\kappa_n^4=n^2+2n$ and $n^2\Prob(A_n^c)\to0$. Finally $\frac12\E[M^2+v]=\frac{\lambda-\nu}\beta+\nu+\frac\nu\beta=\frac\lambda\beta+\nu$, and $\E[ZM]=\beta\E v=2\nu$ by \cref{lem:generalization:parts}(iii).
\end{proof}

\begin{heuristic}[The Bayes training error]
Write $\Lambda_i=\log\bigl(p(X_i\mid w)/q(X_i)\bigr)$, which is $(c-c_0)x_i\varepsilon_i-\frac12(c-c_0)^2x_i^2$. Then $n(T_n(\beta)-S_n)=-\sum_i\log\Ewb{e^{\Lambda_i}}$. The cumulant expansion $\log\Ewb{e^{\Lambda_i}}\approx\Ewb{\Lambda_i}+\frac12\Var^\beta_w(\Lambda_i)$, whose third-order terms total $O(n^{-1/2})$, gives $n(T_n(\beta)-S_n)\approx n(T_{n,\mathrm g}(\beta)-S_n)-\frac12V_n$, and $V_n\approx\Var^\beta_w(c)\sum_ix_i^2\varepsilon_i^2\approx\Var_{\pi_n}(t)$, which tends in law to $v(Z)$. So $n\E[T_n(\beta)-S_n]\to\frac\lambda\beta-\nu-\frac\nu\beta=\frac{\lambda-\nu}\beta-\nu$, and $\frac\beta2\E V_n\to\nu$: our $\nu$ is Watanabe's singular fluctuation in both of his definitions below.
\end{heuristic}

The four limits satisfy two linear relations that do not involve $\lambda$: the generalization errors exceed the training errors by $2\nu$, for the Bayes and for the Gibbs pair, which is $2\beta$ times the gap $\nu/\beta$ between the Gibbs and the Bayes training errors. These are Watanabe's \term{equations of state}\index{equations of state}, named so because they hold for every true distribution, model, prior and singularity; here they are Stein's identity. The sum $2\lambda/\beta$ of the Gibbs limits is the virial identity, the only relation in which the singularity appears. The \term{singular fluctuation}\index{singular fluctuation} $\nu$ is not determined by $\lambda$; \cref{tab:generalization:nu} lists it for prior densities $\propto\abs c^{2\lambda-1}$.

\begin{table}[t]
\centering
\small
\setlength{\tabcolsep}{5pt}
\begin{tabular}{@{}lcccccc@{}}
\toprule
prior density at the truth & $\lambda$ & $\nu(\frac12)$ & $\nu(1)$ & $\nu(2)$ & $\lambda-2\nu(1)$ & $\lambda+\nu(1)$\\
\midrule
$\abs c^{-2/3}$, cubic cusp & $\frac16$ & $0.2263$ & $0.2784$ & $0.3458$ & $-0.3901$ & $0.4450$\\
$\abs c^{-1/2}$, cusp & $\frac14$ & $0.3075$ & $0.3487$ & $0.3967$ & $-0.4474$ & $0.5987$\\
constant, regular line & $\frac12$ & $0.5$ & $0.5$ & $0.5$ & $-0.5$ & $1$\\
$\abs c^{1/2}$ & $\frac34$ & $0.6500$ & $0.6099$ & $0.5727$ & $-0.4699$ & $1.3599$\\
$\abs c^{-1}$ on $\R^2$, variant & $\frac12$ & $0.6157$ & $0.6997$ & $0.7976$ & $-0.8993$ & $1.1997$\\
\bottomrule
\end{tabular}
\caption{The singular fluctuation $\nu(\beta)=\frac\beta2\E[v(Z)]$ of the limit posterior $\propto e^{-\beta(t-Z)^2/2}\abs t^{2\lambda-1}$, and the limits of $n$ times the Bayes training and Gibbs generalization errors at $\beta=1$ (the Bayes generalization limit is $\lambda$, the Gibbs training limit $\lambda-\nu$). The product model has the row of the regular line; the last row is the variant described in the text. Quadrature with $120$ Gauss--Hermite nodes in $Z$, under which the virial and Stein identities hold to $10^{-12}$ (for the variant, $160$ Gauss--Laguerre nodes in $\frac12\abs Z^2$). Script: \texttt{09-generalization-quartet.py}.}\label{tab:generalization:nu}
\end{table}

In the regular row $\nu=\lambda=\frac12$ at every temperature, and at $\beta=1$ the Gibbs error is twice the Bayes error. The product model has the same row, because the logarithm in its prior disappears in the limit: at the level of $\frac1n$ it is a regular one-parameter model, whose singularity shows only through $\lambda<\frac d2$ and the slow corrections of $m=2$. In the cusps $\nu\ne\lambda$: the cubic cusp generalizes at $\frac1{6n}$, but its training error lies $\frac{2\nu}n=\frac{0.557}n$ below, the gap that WAIC estimates (\cref{ch:waic}). And the Bayes error depends on $\beta$ through $\nu$: for the cusp its limits are $0.192$, $0.250$ and $0.323$ at $\beta=\frac12,1,2$ (\cref{ex:generalization:temperature}).

The two integrations by parts also work in $\R^k$, where a weight homogeneous of degree $\alpha>-k$ gives $\lambda=\frac{k+\alpha}2$. The last row of \cref{tab:generalization:nu} is the variant $y=a\,(b^\top x)+\varepsilon$, $x\sim\Normal(0,I_2)$, with $a$ uniform on $[-1,1]$ and $b$ uniform on the unit disk: $ab$ has the density $\frac1\pi(\abs c^{-1}-1)$ on the disk, so $k=2$, $\alpha=-1$, $\lambda=\frac12$, $m=1$, and $\nu>\lambda$. We did not extend \cref{lem:generalization:posterior} to $\R^k$, so for this model the limits are argued, not proved.

Watanabe's theorem states the same structure for analytic models. We quote it with the hypotheses of the paper that proves it.

\begin{theorem}[Watanabe's generalization theorem and the equations of state]\label{thm:generalization:main}
Let the model $p(x\mid w)$, the prior $\varphi$ and the truth $q$ satisfy the following conditions, with $f(x,w)=\log\bigl(q(x)/p(x\mid w)\bigr)$.
\begin{enumerate}[label=(\roman*),nosep]
\item $W$ is compact, the closure of an open set, and $W=\{w:\pi_1(w)\ge0,\dots,\pi_k(w)\ge0\}$ for analytic functions $\pi_j$; $\varphi=\varphi_0\varphi_1$ with $\varphi_0>0$ infinitely differentiable and $\varphi_1\ge0$ analytic.
\item For some $s\ge6$, $w\mapsto f(\cdot,w)$ is an analytic map from an open set $W'\subset\C^d$ containing $W$ to $L^s(q)$.
\item The truth is realizable, $W_0=\{w:p(\cdot\mid w)=q\}\ne\emptyset$; on an open set $W^*\subset\C^d$ containing $W$ the envelope $\bar f(x)=\sup_{W^*}\abs{f(x,w)}$ has $\E_X[\bar f(X)^s]<\infty$, and $\int\bar f(x)^2\sup_{K(w)\le\epsilon}p(x\mid w)\dd x<\infty$ for some $\epsilon>0$.
\end{enumerate}
Let $\beta>0$, let $\lambda$ be the learning coefficient, and write $B_{\mathrm g}=G_n(\beta)-S$, $B_{\mathrm t}=T_n(\beta)-S_n$, $G_{\mathrm g}=G_{n,\mathrm g}(\beta)-S$, $G_{\mathrm t}=T_{n,\mathrm g}(\beta)-S_n$. Then:
\begin{enumerate}[label=(\alph*),nosep]
\item $nB_{\mathrm g}$, $nB_{\mathrm t}$, $nG_{\mathrm g}$, $nG_{\mathrm t}$ converge in law to random variables $B_{\mathrm g}^*$, $B_{\mathrm t}^*$, $G_{\mathrm g}^*$, $G_{\mathrm t}^*$, and their expectations converge to those of the limits;
\item (the equations of state) $\E B_{\mathrm g}^*-\E B_{\mathrm t}^*=2\beta(\E G_{\mathrm t}^*-\E B_{\mathrm t}^*)$ and $\E G_{\mathrm g}^*-\E G_{\mathrm t}^*=2\beta(\E G_{\mathrm t}^*-\E B_{\mathrm t}^*)$;
\item $nG_{\mathrm g}+nG_{\mathrm t}\to2\lambda/\beta$ in probability, so $\E G_{\mathrm g}^*+\E G_{\mathrm t}^*=2\lambda/\beta$;
\item with the singular fluctuation $\nu=\nu(\beta)=\beta(\E G_{\mathrm t}^*-\E B_{\mathrm t}^*)$,
\[
\E B_{\mathrm g}^*=\frac{\lambda-\nu}\beta+\nu,\quad\E B_{\mathrm t}^*=\frac{\lambda-\nu}\beta-\nu,\quad
\E G_{\mathrm g}^*=\frac\lambda\beta+\nu,\quad\E G_{\mathrm t}^*=\frac\lambda\beta-\nu .
\]
\end{enumerate}
In particular, at $\beta=1$, $\E[G_n]-S=\frac\lambda n+o(\frac1n)$.
\end{theorem}

This is Theorem~1, Theorem~2, Theorem~3 and Corollary~3 of \textcite{Watanabe2010a}, numbered as in its arXiv version, whose conditions (A.1)--(A.3) are our (i)--(iii); the value $\E B_{\mathrm g}^*=\lambda$ at $\beta=1$ ends his Corollary~2, which he says was first found in his papers of 1999 and 2001 \parencite{Watanabe1999,Watanabe2001a}. \Textcite[Lemma~3]{Watanabe2010} states the limits of (d) for the Bayes errors without realizability, with the proof for unrealizable singular truths in \textcite{Watanabe2010d}, under relatively finite variance (\cref{thm:freeenergy:main}), essential uniqueness and $s\ge8$, with $L(w_0)$, $L_n(w_0)$ in place of $S$, $S_n$ and $\nu=\lim_n\frac\beta2\E[V_n]$; at $\beta=1$ this is $\E[G_n]=L(w_0)+\frac\lambda n+o(\frac1n)$ and $\E[T_n]=L(w_0)+\frac{\lambda-2\nu}n+o(\frac1n)$ \parencite{Watanabe2024}.

\begin{proofsketch}
The plan is the proof of \cref{thm:generalization:onecoef} on a parameter space on which $K$ is a monomial, with the Gaussian observation $Z$ replaced by a Gaussian process. As in the proof of \cref{thm:freeenergy:main}, the posterior concentrates on $\{K<\epsilon\}$, Hironaka's theorem (\cref{thm:toolkit:hironaka}) gives local coordinates with $K=u^{2k}$ and $f=u^ka(x,u)$, and with $t=nu^{2k}$ the energy is $nK_n=t-\sqrt t\,\xi_n(u)$ up to the sign of $u^k$, where $\xi_n$ converges in law to a Gaussian process $\xi$ (\cref{thm:toolkit:empirical}). To leading order the posterior becomes the measure $t^{\lambda-1}e^{-\beta t+\beta\sqrt t\,\xi(u)}\dd t\,\dd u^*$, the analogue of $\pi_Z$, with averages $\langle\cdot\rangle$, and a cumulant expansion writes $n$ times each error as such an average up to $o_p(1)$ \parencite[\S5]{Watanabe2024}. The integration by parts $\int_0^\infty\partial_t(t^\lambda e^{-\beta t+\beta\sqrt t\,\xi})\dd t=0$ is the virial identity $\beta\langle t\rangle=\lambda+\frac\beta2\langle\sqrt t\,\xi\rangle$, which gives (c) since $nG_{\mathrm g}\approx\langle t\rangle$ and $nG_{\mathrm t}\approx\langle t-\sqrt t\,\xi\rangle$. Watanabe represents $\xi$ as the limit in law of $n^{-1/2}\sum_ia(X_i,u)g_i$ with independent standard Gaussian $g_i$, and Stein's identity in each $g_i$ gives (b); he writes that Theorems 2 and 3 are ``essentially proved by partial integration'' \parencite[\S5]{Watanabe2010a}. Then (d) is linear algebra, and moment bounds on $\sup_u\abs{\xi_n}$ give (a). Omitted are the resolution, the empirical processes on the resolved space and the moment bounds: \S4 and the appendix of \textcite{Watanabe2010a}, and Chapters~5 and~6 of \textcite{Watanabe2009}.
\end{proofsketch}

In a one-coefficient model the sketch's $t$ is our $\frac12t^2$ to leading order and its $\sqrt t\,\xi$ is our $tZ$. The free energy of \cref{ch:freeenergy} gives $\lambda$ by a different road, and predicts the slow convergence measured in \cref{sec:generalization:measuring}.

\begin{heuristic}[From the free energy to the generalization error]
Since $F_{n+1}-F_n=-\log p^*(X_{n+1})$, the mean over $X_{n+1}$ gives the exact identity $\E[G_n]=\E[F_{n+1}]-\E[F_n]$. Differencing the expansion $\E[F_n]-nS=\lambda\log n-(m-1)\log\log n+c+o(1)$ \parencite{Watanabe2024} term by term, which needs an unchecked regularity condition, gives $\E[G_n]-S\approx\frac\lambda n-\frac{m-1}{n\log n}$. In one-coefficient models a second argument agrees: to first order in $1/\log\kappa_n$ the rescaled weight $\abs t^\alpha(1-\log\abs t/\log\kappa_n)^{m-1}$ is $\abs t^{\alpha'}$ with $\alpha'=\alpha-(m-1)/\log\kappa_n$, and \cref{lem:generalization:parts}(iv) gives $\frac12(1+\alpha')\approx\lambda-\frac{m-1}{\log n}$.
\end{heuristic}

\section{Measuring it}\label{sec:generalization:measuring}

The script \texttt{09-generalization-onecoef.py} computes the Bayes generalization error of the regular line, the product model, the cusp and the cubic cusp, data set by data set, without sampling. By \cref{lem:generalization:reduction} a data set enters only through the independent $\kappa_n^2\sim\chi^2_n$ and $Z_n\sim\Normal(0,1)$, so the script draws $1000$ independent data sets through these two statistics for each model and each $n\in\{10^2,10^3,\dots,10^6\}$ (seed $900+10j+\log_{10}n$, $j=0,\dots,3$ in the order just listed).\sidenote{Drawing the two statistics instead of $n$ pairs changes nothing in law, and at $n=10^6$ saves $2\times10^9$ Gaussian draws per model.} It computes $G_n-S=\E_x[D(\mu_x)]$ by $16$-point Gauss--Hermite rules in $x$ and $\varepsilon$, and composite Gauss--Legendre rules in $t$ after the substitution $t=\pm u^p$, $p=1,4,2,3$, which removes the singularity of $\rho$ at $0$. Tripling the rules changes $n(G_n-S)$ by less than $3\times10^{-9}$; on the regular line, with the uniform prior on $[-1,1]$, the result matches the flat-prior closed form of \cref{comp:generalization:line} to within $1.8\times10^{-7}$ on each of its $5000$ data sets. The run takes about a minute.

\begin{figure}[t]
\centering
\begin{tikzpicture}
\begin{groupplot}[group style={group size=2 by 2,horizontal sep=12mm,vertical sep=13mm},
  splot,width=0.5\linewidth,height=0.33\linewidth,xmode=log,xmin=50,xmax=2e6,ymin=0,ymax=1.12,
  xtick={1e2,1e3,1e4,1e5,1e6},ytick={0,0.25,0.5,0.75,1},yticklabels={$0$,$0.25$,$0.5$,$0.75$,$1$},
  title style={font=\small,yshift=-0.6ex}]
\nextgroupplot[title={regular line $y=ax$},ylabel={$n(\E[G_n]-S)$}]
\addplot[regularpred,domain=50:2e6,samples=2] {0.5};
\addplot[theory,domain=50:2e6,samples=2] {0.5};
\addplot[theory2] table[x=n,y=line] {figures/data/09-generalization-pred.dat};
\addplot[measured,ebars] table[x=n,y=line,y error plus=lineci,y error minus=lineci] {figures/data/09-generalization-onecoef.dat};
\nextgroupplot[title={product model $y=abx$}]
\addplot[regularpred,domain=50:2e6,samples=2] {1};
\addplot[theory,domain=50:2e6,samples=2] {0.5};
\addplot[theory2] table[x=n,y=product] {figures/data/09-generalization-pred.dat};
\addplot[measured,ebars] table[x=n,y=product,y error plus=productci,y error minus=productci] {figures/data/09-generalization-onecoef.dat};
\nextgroupplot[title={cusp $y=ab^2x$},ylabel={$n(\E[G_n]-S)$},xlabel={sample size $n$}]
\addplot[regularpred,domain=50:2e6,samples=2] {1};
\addplot[theory,domain=50:2e6,samples=2] {0.25};
\addplot[theory2] table[x=n,y=cusp] {figures/data/09-generalization-pred.dat};
\addplot[measured,ebars] table[x=n,y=cusp,y error plus=cuspci,y error minus=cuspci] {figures/data/09-generalization-onecoef.dat};
\nextgroupplot[title={cubic cusp $y=ab^3x$},xlabel={sample size $n$},
  legend style={font=\scriptsize,at={(0.97,0.55)},anchor=east}]
\addplot[regularpred,domain=50:2e6,samples=2] {1};
\addlegendentry{regular: $\frac d2$}
\addplot[theory,domain=50:2e6,samples=2] {1/6};
\addlegendentry{singular: $\lambda$}
\addplot[theory2] table[x=n,y=cubic] {figures/data/09-generalization-pred.dat};
\addlegendentry{finite $n$: $\frac12\E\,\bar t^{\,2}$}
\addplot[measured,ebars] table[x=n,y=cubic,y error plus=cubicci,y error minus=cubicci] {figures/data/09-generalization-onecoef.dat};
\addlegendentry{measured}
\end{groupplot}
\end{tikzpicture}
\caption{The Bayes generalization error, times $n$, of four one-coefficient models at the truth $0$; inputs $\Normal(0,1)$, unit noise, uniform prior on $[-1,1]^d$. Dots: means over $1000$ independent data sets at each $n$, each by exact quadrature, with Student-$t$ $95\%$ intervals (seeds $900+10j+\log_{10}n$). Bronze: the limit $\lambda$ of \cref{thm:generalization:onecoef}. Dotted bronze: the finite-$n$ prediction $\frac12\E[\bar t(\sqrt n,Z)^2]$, half the mean square of the exact posterior mean of $t$ at $\kappa_n=\sqrt n$ (for the regular line, the closed form of \cref{comp:generalization:line}). Slate dashed: the regular prediction $\frac d2$, equal to $\lambda$ for the regular line. Script: \texttt{09-generalization-onecoef.py}.}\label{fig:generalization:measured}
\end{figure}

\Cref{fig:generalization:measured} shows the result. None of the three two-parameter models comes near the regular prediction $1$: at $n=10^6$ the cubic cusp measures $0.164\pm0.024$ against $\lambda=\frac16$ and the cusp $0.244\pm0.031$ against $\frac14$ (the product model is discussed below), while the one-parameter regular line measures $0.504\pm0.047$ against $\frac12$. Nineteen of the twenty intervals contain the dotted prediction; the twentieth, the regular line at $n=10^4$ with $0.561\pm0.049$ against the exact $0.50003$, is $2.5$ standard errors high, as one interval in twenty may be. Half the squared posterior mean accounts for nearly all of the error: its mean differs from the exact error by at most $0.009$ at $n=100$, $0.001$ at $n=1000$, and $10^{-4}$ from $n=10^4$ on, as an $O(n^{-2})$ remainder should.

The singular models approach $\lambda$ slowly. The product model has $m=2$: its finite-$n$ prediction is $0.302$, $0.365$, $0.397$, $0.417$ and $0.430$ at $n=10^2,\dots,10^6$, close to the heuristic $\frac12-\frac1{\log n}$ ($0.283$, $0.355$, $0.391$, $0.413$, $0.428$), and the measured values $0.277\pm0.029$, $0.363\pm0.038$, $0.436\pm0.043$, $0.433\pm0.041$ and $0.422\pm0.038$ follow it; the prediction is still $0.4785$ at $n=10^{20}$ and $0.4957$ at $n=10^{100}$.\sidenote{At $n=10^{100}$ the product model is still $0.9\%$ below its limit, as $\frac12-\frac1{\log n}$ predicts.} The cusps have $m=1$ but a second term in their prior densities, the $-1$ in $\frac12(\abs c^{-1/2}-1)$, which multiplies the rescaled weight by $1-\abs t^{1/2}\kappa_n^{-1/2}$, and by $1-\abs t^{2/3}\kappa_n^{-2/3}$ for the cubic cusp. Their predictions approach $\lambda$ like $n^{-1/4}$ and $n^{-1/3}$: $\frac14-0.2303=0.0197$ at $n=10^4$ and $0.0059$ at $10^6$, and $\frac16-0.1648=0.0019$ at $10^6$. The second term is the next pole of the zeta function (\cref{ch:volume}), at $-\frac12$, and the rate is the distance between the poles; we know of no theorem that says this in general. A single data set's error does not concentrate: at $n=10^6$ its standard deviation is $0.75$, $0.51$ and $0.38$ for the regular line, the cusp and the cubic cusp, against $0.71$, $0.51$ and $0.42$ for the limit law (\cref{ex:generalization:law}).

\section{The ledger}\label{sec:generalization:ledger}

The expected generalization error of the Bayes predictive distribution is $\lambda/n$, not $\frac d{2n}$. For one-coefficient regressions with Gaussian inputs and noise this is proved here completely, for every learning coefficient and multiplicity: the divergence to the predictive is half the squared posterior mean, the rescaled posterior tends to a Gaussian location posterior with a homogeneous prior weight, and two integrations by parts show that the mean square of its mean is $2\lambda$. At other temperatures the same proof gives $(\frac{\lambda-\nu}\beta+\nu)\frac1n$, and the Gibbs errors take the values of Watanabe's equations of state, which in these models are Stein's identity (the Bayes training error is argued); Watanabe's theorem proves all four for analytic models. The cubic cusp generalizes three times better than one regular parameter; the product model, at the level of $\frac1n$, exactly as well as one.

\ledgercard{the Bayes generalization error $\E[G_n]-S$, the expected divergence from the truth to the predictive}{$\frac d{2n}+o(\frac1n)$ (\cref{thm:regular:aic}); training error $-\frac d{2n}$; Gibbs error $\frac dn$}{$\frac\lambda n+o(\frac1n)$ at $\beta=1$: proved for one-coefficient models (\cref{thm:generalization:onecoef}), Watanabe's theorem in general (\cref{thm:generalization:main}); $(\frac{\lambda-\nu}\beta+\nu)\frac1n$ at temperature $\beta$; training error $(\lambda-2\nu)\frac1n$}{the cubic cusp $y=ab^3x$: $\frac1{6n}$ against $\frac1n$ for $d=2$; the error is half the mean square of the posterior mean $M(Z)$, and $\E M^2=\frac13$; $\nu(1)=0.278$}{$n(\E[G_n]-S)$ at $n=10^6$: $0.164\pm0.024$ (cubic cusp), $0.244\pm0.031$ (cusp), $0.504\pm0.047$ (regular line), $0.422\pm0.038$ (product model, $m=2$, approaching $\frac12$ like $\frac12-\frac1{\log n}$)}{the exact posterior average at $\beta=1$; a fixed realizable truth as $n\to\infty$; one coefficient, Gaussian inputs and noise (the proof); analyticity, compactness, and realizability or relatively finite variance (Watanabe)}

\section{Where the theory stops}\label{sec:generalization:stops}

\paragraph{What was assumed.} The theorems describe the exact posterior average at $\beta=1$. A point estimate does not inherit $\lambda/n$: for the cubic cusp the maximum-likelihood estimate, also the maximum a posteriori estimate under the uniform prior on $(a,b)$, has error $\frac1{2n}$, three times the Bayes error (\cref{ex:generalization:mle}), and for analytic models the maximum-likelihood error is governed by the supremum of a Gaussian process over the true parameters \parencite{Watanabe2024}. The truth is fixed as $n$ grows: at a truth $c_0\ne0$ of the cubic cusp, $\lambda=\frac12$ and the gain disappears. For a truth at distance $\delta$ from the singularity, the limit posterior with its truth moved to $\sqrt n\,\delta$ gives an error of about $\frac1{6n}$ while $n\delta^2$ is small and $\frac1{2n}$ once it is large, but $0.77/n$ at $\sqrt n\,\delta=2.8$: near the singularity the bet is lost, in a crossover like those of \cref{ch:phases}. Our proof used one coefficient and Gaussian inputs and noise; Watanabe's theorem needs analyticity, which ReLU networks lack, a compact parameter set, moments of order $6$, and a prior of the form $\varphi_0\varphi_1$: with a smooth but non-analytic prior factor $e^{-1/\norm w^2}$ at the only true parameter the errors need not be proportional to $\frac1n$ \parencite[\S2]{Watanabe2010a}. And the theorem is about the mean over data sets: under the limit law of \cref{thm:generalization:onecoef}(ii), a data set of the cubic cusp still has an error above $\frac1{2n}$ with probability $0.084$.

\paragraph{Where it breaks.} The $o(\frac1n)$ term can be large at every practical sample size: the product model's prediction is $0.430$ at $n=10^6$ against its limit $\frac12$, and a second term in the prior slows the cusps to $n^{-1/4}$ and $n^{-1/3}$. In unrealizable pairs without relatively finite variance the theorem fails: for Watanabe's non-renormalizable example, whose free energy \cref{ch:freeenergy} measures, he argues that the learning curve is of order $n^{-2/3}$ \parencite{Watanabe2010d}. The theorem says nothing when $d$ is comparable to $n$, and trained networks are not posterior averages (\cref{ch:sgd}). In Bayesian deep learning, \textcite{WenzelEtAl2020} found carefully sampled posterior predictives worse than SGD point estimates and improved by cold posteriors, $\beta>1$,\sidenote{Cold means a temperature $T=1/\beta$ below $1$: the likelihood is raised to the power $\beta>1$, as though each observation were counted $\beta$ times.} while \textcite{IzmailovEtAl2021} found tempering unnecessary and attribute the effect largely to data augmentation. Even in a one-coefficient model the best temperature depends on where the prior puts its mass (\cref{prop:generalization:quartet}); whether $\nu$ has anything to do with the cold-posterior effect has, as far as we know, not been tested.

\paragraph{Open problems.} Our computations suggest a remainder $-\frac{m-1}{n\log n}$ when $m\ge2$, and of order $n^{-1-(\lambda_2-\lambda)}$ when the next pole of the zeta function is at $-\lambda_2$; we know of no general theorem on the rate of the $o(\frac1n)$ term. The singular fluctuation is not determined by $\lambda$: Watanabe conjectured that it is governed by a generalized Fisher information on the set of true parameters \parencite[\S5]{Watanabe2010a}, and we know of no closed form for $\nu$ even for the power-law class of \cref{tab:generalization:nu}. Non-asymptotic bounds on $\E[G_n]-S$ for singular models, and the learning curve when $d$ grows with $n$, are open.

\begin{exercises}
\item \consolidation (The fourth order.) Continue \cref{comp:generalization:divergence} to degree $4$: with $\bar u$ and $\gamma_0$ the posterior mean and variance of $u$, show that $D_x=\frac12\bar u^2+\frac14\gamma_0^2-\frac12\bar u^2\gamma_0+O(u^6)$. Check it against \cref{comp:generalization:line} expanded in $v$, and show that its mean, $-\frac3{4n^2}+O(n^{-3})$, is the difference between the exact $\frac1{2n}+\frac1{4n^2}$ and $\frac1{2(n-2)}$. \hint{the terms of degree $4$ come from $\frac12\E A^2$, $\frac13\E A^3$ and $\frac14\E A^4$; $\E[\varepsilon^2\mathrm{He}_2(\varepsilon)]=2$.}\label{ex:generalization:fourth}

\item \consolidation (The Riccati equation.) For the weight $\abs t^\alpha$ at $\beta=1$ let $r=h'/h$, so that $M=z+r$. Show that $r'=\alpha-zr-r^2$ (\cref{thm:toolkit:stein}(iv)) is \cref{lem:generalization:parts}(i) rewritten with $M'=v$, and derive $\E[M(Z)^2]=1+\alpha$ by applying Stein's identity to $r$.\label{ex:generalization:riccati}

\item \consolidation (Gibbs and Bayes on the regular line.) The limit posterior of the regular line at inverse temperature $\beta$ is $\Normal(Z,1/\beta)$. Compute the four limits of \cref{thm:generalization:main}(d) directly, $\frac12$, $-\frac12$, $\frac12+\frac1{2\beta}$ and $\frac1{2\beta}-\frac12$, and check the equations of state with $\nu=\frac12$.\label{ex:generalization:regular}

\item \consolidation (One data set.) From \cref{thm:generalization:onecoef}(ii), show that the limit law of $n(G_n-S)$ has standard deviation $\frac1{\sqrt2}$ for the regular line ($0.423$ for the cubic cusp, numerically), and that the limit probability of $n(G_n-S)>\frac12$ is $2\Prob(Z>z_*)$ with $M(z_*)=1$: $0.084$ for the cubic cusp, where $z_*=1.730$, against $0.317$ for the regular line.\label{ex:generalization:law}

\item \byhand (A one-sided coefficient.) For $y=a^2x$ with $a$ uniform on $[-1,1]$ and the truth $0$, show that $c=a^2$ has the density $\frac1{2\sqrt c}$ on $(0,1]$, that $\lambda=\frac14$, and that the virial identity holds for the one-sided weight $t^{-1/2}\ind\{t>0\}$, so that $\E[M(Z)^2]=\frac12$. \hard{} Adapt the denominator bound of \cref{lem:generalization:posterior}, which must allow a factor $e^{-z^2/2}$ for $z<0$, and prove $n(\E[G_n]-S)\to\frac14$. \hint{for $z<0$, integrate over $0<t<1/(1+\abs z)$.}\label{ex:generalization:onesided}

\item \byhand (Maximum likelihood against Bayes.) For the cubic cusp at the truth $0$, show that the maximum-likelihood estimate of $c=ab^3$, clipped to $[-1,1]$, is also the maximum a posteriori estimate under the uniform prior on $(a,b)$, and that $n\E[\KL(q\,\|\,p(\cdot\mid\hat w))]\to\frac12$. For which $\lambda$ in the power-law class is the Bayes predictive worse?\label{ex:generalization:mle}

\item \byhand (The chain.) For $y=a_1a_2a_3x$ with the uniform prior on $[-1,1]^3$ and the truth $0$, $\rho(c)=\frac14(\log\frac1{\abs c})^2$. Show that the effective-exponent argument gives $n(\E[G_n]-S)\approx\frac12-\frac2{\log n}$, and compare with the finite-$n$ predictions $0.303$, $0.363$, $0.396$ at $n=10^4,10^6,10^8$.\label{ex:generalization:chain}

\item \variation (Temperature.) For the weight $\abs t^\alpha$, show that $\pi_z$ at inverse temperature $\beta$ is the law of $\beta^{-1/2}t'$, with $t'$ from the limit posterior at $\beta=1$ given the observation $\sqrt\beta z$. Deduce $\lim_nn\E[G_n(\beta)-S]=\frac1{2\beta}\E[M_1(\sqrt\beta Z)^2]$, which tends to $2\lambda^2$ as $\beta\to0$ and to $\frac12$ as $\beta\to\infty$. Which temperature is best for the cusp at its singular truth, and which for $\lambda=\frac34$?\label{ex:generalization:temperature}

\item \variation (An asymmetric prior.) Let $\ell_+\ne\ell_-$ in \cref{ass:generalization:onecoef}. Explain why the limit of \cref{thm:generalization:onecoef} does not depend on $\ell_+/\ell_-$ while $\nu$ does; compute $\nu(1)$ for $\lambda=\frac14$ and $\ell_-/\ell_+=0.2$ (answer: $0.3065$, against $0.3487$).\label{ex:generalization:asymmetric}

\item \measurement (The training error.) Extend the script of \cref{sec:generalization:measuring} to the Bayes training error of the cubic cusp, on simulated data sets of size $n=10^3$ and $10^4$ (the training points are now needed). Compare $n\E[T_n-S_n]$ with $\lambda-2\nu=-0.390$ and $\frac12\E[V_n]$ with $\nu=0.278$.\label{ex:generalization:training}

\item \research\hard (Rates.) Show that for the cusp $n(\E[G_n]-S)-\frac14$ is of exact order $n^{-1/4}$. Conjecture the rate in terms of the first two poles of the zeta function and their orders, test it on the cubic cusp, the chain, and the density $\frac3{16}\abs c^{-1/2}(1+\abs c)$ on $[-1,1]$, whose second pole is at $-\frac34$, and find what the literature proves.\label{ex:generalization:rates}
\end{exercises}

\notesandrefs
For regular models the learning curve $\frac d{2n}$ is classical: with the training error $-\frac d{2n}$ it makes up Akaike's penalty $\frac dn$ \parencite{Akaike1974}, and \textcite{AmariMurata1993} proved the entropic errors $H_0\pm\frac m{2t}$ of a faithful regular model with $m$ parameters and $t$ examples for maximum likelihood, with similar relations for Bayes and Gibbs learning. Summed over $n$, the learning curve is the redundancy of the Bayes mixture, $\frac d2\log n$ plus an explicit constant \parencite{ClarkeBarron1990}.

For singular models Watanabe obtained $\lambda/n$ first as the increase of the free energy \parencite{Watanabe1999,Watanabe2001a}. The equations of state, posted in 2007 and published as \textcite{Watanabe2010a}, prove it directly; the 2009 book counts them among its four main formulas \parencite{Watanabe2009}. \Textcite{Watanabe2010} states the Bayes limits for unrealizable truths and proves WAIC asymptotically equivalent to cross-validation (\cref{ch:waic}); \textcite{Watanabe2024} reviews the theory at $\beta=1$, and the 2018 textbook \parencite{Watanabe2018} develops it with exercises.

Tweedie's formula was first published by \textcite{Robbins1956}, who credits personal correspondence with Maurice Kenneth Tweedie \parencite{Efron2011}; Efron bases empirical Bayes corrections for selection bias on it and gives the second-order form, our $M'=\beta v$. The proofs of \cref{thm:generalization:onecoef} and \cref{prop:generalization:quartet} are the book's own. The local learning coefficient of \textcite{LauEtAl2025} estimates $\lambda$ near a trained network (\cref{ch:llc}); its link to that network's test loss runs through a posterior average that SGD does not compute (\cref{ch:sgd}).

\paragraph{Translating the literature.} \Textcite{Watanabe2010a} writes $E_w$ for our $\Ewb{\cdot}$; his Bayes and Gibbs errors $B_g$, $B_t$, $G_g$, $G_t$ are our $G_n(\beta)-S$, $T_n(\beta)-S_n$, $G_{n,\mathrm g}(\beta)-S$, $T_{n,\mathrm g}(\beta)-S_n$, his $B_g^*$ is the limit in law of $nB_g$, and his losses $BL_g$, $BL_t$, $GL_g$, $GL_t$ are our $G_n(\beta)$, $T_n(\beta)$, $G_{n,\mathrm g}(\beta)$, $T_{n,\mathrm g}(\beta)$. \Textcite{Watanabe2010} writes $BgL(n)$ and $BtL(n)$ for our $G_n(\beta)$ and $T_n(\beta)$, $Bg(n)=BgL(n)-L(w_0)$, and $V(n)$ for our $V_n$. \Textcite{Watanabe2024} works at $\beta=1$ with $\theta$ and $\pi(\theta)$ for our $w$ and $\varphi(w)$, and $\langle\cdot\rangle$ for the renormalized posterior. \Textcite{Efron2011} writes $z\mid\mu\sim\Normal(\mu,\sigma^2)$ and $l=\log f$ for the log marginal density: his $z$, $\mu$, $\sigma^2$ and $f$ are our $Z$, $t$, $1/\beta$ and $h$ up to a constant factor.

\printchapterbib
\end{refsection}
